\documentclass[11pt,a4paper]{article}

\usepackage[T1]{fontenc}
\usepackage[utf8]{inputenc}
\usepackage[ngerman]{babel}
\usepackage{amsmath}
\usepackage{amssymb}
\usepackage{mathtools}
\usepackage{booktabs}
\usepackage{geometry}
\usepackage[hidelinks]{hyperref}

\geometry{left=28mm,right=28mm,top=25mm,bottom=28mm}

\newcommand{\R}{\mathbb{R}}
\newcommand{\Vol}{\operatorname{Vol}}
\newcommand{\conv}{\operatorname{conv}}

\title{Exakte Volumenberechnung im\\Optimal-Trade-off-Query-Algorithmus}
\author{}
\date{September 2026}

\begin{document}

\maketitle

\begin{abstract}
Dieses Dokument beschreibt die geometrisch exakte Bestimmung von
Kandidaten- und Query-Antwortwahrscheinlichkeiten im
Optimal-Trade-off-Query-Algorithmus. Der zulässige Gewichtsraum wird als
konvexes Polytop dargestellt. Beantwortete Queries und
Optimalitätsbedingungen schneiden dieses Polytop durch lineare Halbräume.
Nach der Elimination der Gewichtssummengleichung werden die Ecken der
entstehenden Polytopen bestimmt und ihre intrinsischen Volumina berechnet.
Die gesuchten Wahrscheinlichkeiten ergeben sich anschließend als
Volumenquotienten. Die Wahrscheinlichkeiten dienen ausschließlich zur
Gewichtung der Zweige der bereits vorhandenen Zielfunktion; es wird kein
zusätzliches Entropie- oder Regret-Ziel eingeführt.
\end{abstract}

\tableofcontents

\section{Ausgangsproblem und Notation}

Es seien
\begin{itemize}
  \item $m$ Ziele,
  \item $n$ Handlungsalternativen und
  \item $x_k\in\R^m$ der Zielwertvektor der Alternative $k$
\end{itemize}
gegeben. Der unbekannte Gewichtsvektor wird mit
\[
  w=(w_1,\dots,w_m)^\top\in\R^m
\]
bezeichnet. Für jedes Gewicht gelten
\[
  w_i\geq 0
  \qquad\text{für alle }i\in\{1,\dots,m\}
\]
und
\[
  \sum_{i=1}^{m}w_i=1.
\]

Der aggregierte Nutzen der Alternative $k$ ist
\[
  U_k(w)=x_k^\top w.
\]
Eine Alternative wird als optimal bezeichnet, wenn ihr aggregierter Nutzen
mindestens so groß ist wie der Nutzen jeder anderen Alternative.

\section{Der Gewichtsraum als konvexes Polytop}

Ohne beantwortete Queries ist der zulässige Gewichtsraum der Standardsimplex
\[
  W_0
  =
  \left\{
    w\in\R^m
    \;\middle|\;
    w_i\geq 0\text{ für alle }i,
    \ \boldsymbol{1}^\top w=1
  \right\}.
\]
Für drei Ziele ist $W_0$ ein Dreieck mit den Ecken
\[
  (1,0,0)^\top,
  \qquad
  (0,1,0)^\top,
  \qquad
  (0,0,1)^\top.
\]

Allgemein ist $W_0$ ein Simplex. Obwohl $w$ aus $m$ Komponenten besteht,
besitzt $W_0$ aufgrund der Gleichung $\boldsymbol{1}^\top w=1$ nur die
intrinsische Dimension
\[
  d=m-1.
\]
Bei sieben Zielen wird daher ein sechsdimensionales Volumen berechnet.
Das gewöhnliche siebendimensionale Volumen in $\R^7$ wäre null und ist nicht
die gesuchte geometrische Größe.

\subsection{Query-Antworten als lineare Schnitte}

Eine Trade-off-Query hat die Form
\[
  w_a\mathrel{?}s\,w_b,
  \qquad s\geq 0.
\]
Die Antwort ``kleiner'' erzeugt die Ungleichung
\[
  w_a-s w_b\leq 0,
\]
während die Antwort ``größer'' durch
\[
  -w_a+s w_b\leq 0
\]
dargestellt wird. Eine exakte Gleichheitsantwort erzeugt die Gleichung
\[
  w_a-s w_b=0.
\]

Für eine Menge $T$ bereits beantworteter Queries lässt sich der aktuelle
Gewichtsraum somit in der allgemeinen Form
\[
  W(T)
  =
  \left\{
    w\in\R^m
    \;\middle|\;
    Aw\leq b,
    \ Cw=d
  \right\}
\]
schreiben. Alle Nebenbedingungen sind linear. Folglich ist $W(T)$ ein
konvexes Polytop, solange es nicht leer und beschränkt ist.

\section{Optimalitätsregionen der Alternativen}

Alternative $k$ ist für ein Gewicht $w$ optimal, wenn
\[
  U_k(w)\geq U_j(w)
  \qquad\text{für alle }j\neq k.
\]
Mit $U_k(w)=x_k^\top w$ folgt
\[
  x_k^\top w\geq x_j^\top w
\]
beziehungsweise
\[
  (x_j-x_k)^\top w\leq 0.
\]
Die Optimalitätsregion von Alternative $k$ im aktuellen Gewichtsraum ist
damit
\[
  W_k(T)
  =
  W(T)
  \cap
  \bigcap_{j\neq k}
  \left\{
    w\in\R^m
    \;\middle|\;
    (x_j-x_k)^\top w\leq 0
  \right\}.
\]

Auch $W_k(T)$ ist ein konvexes Polytop. Ist die Region leer, kann Alternative
$k$ unter keinem zulässigen Gewichtsvektor optimal sein. Eine Region, die nur
aus einer Fläche niedrigerer Dimension besteht, besitzt bezüglich des
intrinsischen Volumens von $W(T)$ ebenfalls Volumen null.

Die Kandidatenregionen überdecken den Gewichtsraum:
\[
  W(T)=\bigcup_{k=1}^{n}W_k(T).
\]
Verschiedene Regionen können sich auf Grenzen überschneiden, auf denen zwei
oder mehr Alternativen den gleichen Nutzen besitzen. Solche Grenzen haben im
voll-dimensionalen Gewichtsraum Volumen null. Daher gilt
\[
  \Vol(W(T))
  =
  \sum_{k=1}^{n}\Vol(W_k(T)).
\]
Diese Identität wird in den Benchmarks als numerische Konsistenzprüfung
verwendet.

\section{Wahrscheinlichkeiten als Volumenquotienten}

Es wird angenommen, dass der unbekannte Gewichtsvektor innerhalb des aktuell
zulässigen Gewichtsraums gleichverteilt ist. Unter dieser Modellannahme ist
die Wahrscheinlichkeit, dass Alternative $k$ optimal ist,
\[
  p_k(T)
  =
  \Pr(k\text{ ist optimal}\mid T)
  =
  \frac{\Vol(W_k(T))}{\Vol(W(T))}.
\]
Die Wahrscheinlichkeiten erfüllen
\[
  p_k(T)\geq 0
  \qquad\text{und}\qquad
  \sum_{k=1}^{n}p_k(T)=1,
\]
bis auf unvermeidbare Rundungsfehler der Fließkommaarithmetik.

Wichtig ist die Trennung zwischen Geometrie und Entscheidungsziel:
Die Volumenquotienten definieren keine neue Zielfunktion. Sie liefern lediglich
die Wahrscheinlichkeiten, die zur Auswertung der vorhandenen Zielfunktion
benötigt werden.

\section{Berechnung des intrinsischen Volumens}

\subsection{Elimination der Gleichungen}

Algorithmen zur Ecken- und Volumenberechnung benötigen eine Darstellung im
tatsächlichen affinen Raum des Polytops. Sei
\[
  Cw=d
\]
das Gleichungssystem. Zunächst wird ein spezieller Punkt $w_0$ mit
\[
  Cw_0=d
\]
bestimmt. Außerdem sei
\[
  B\in\R^{m\times d_W}
\]
eine orthonormale Basis des Nullraums von $C$. Es gilt also
\[
  CB=0
  \qquad\text{und}\qquad
  B^\top B=I.
\]
Jeder Punkt im affinen Gleichungsraum besitzt nun die eindeutige Darstellung
\[
  w=w_0+Bz,
  \qquad z\in\R^{d_W}.
\]

Setzt man diese Darstellung in die Ungleichungen $Aw\leq b$ ein, erhält man
\begin{align*}
  A(w_0+Bz)&\leq b,\\
  ABz&\leq b-Aw_0.
\end{align*}
Das reduzierte Polytop lautet somit
\[
  P
  =
  \left\{
    z\in\R^{d_W}
    \;\middle|\;
    \widetilde A z\leq\widetilde b
  \right\},
\]
wobei
\[
  \widetilde A=AB
  \qquad\text{und}\qquad
  \widetilde b=b-Aw_0.
\]
Da $B$ orthonormal ist, bleiben Längen und intrinsische Volumina bei dieser
Koordinatentransformation erhalten.

\subsection{Relativer Chebyshev-Mittelpunkt}

Die Halfspace-Intersection-Methode benötigt einen Punkt, der strikt im Inneren
aller nichtredundanten Ungleichungen liegt. Ein beliebiger Machbarkeitspunkt
eines LP-Solvers genügt nicht, weil dieser häufig auf einer Ecke oder Kante
liegt.

Deshalb wird ein relativer Chebyshev-Mittelpunkt berechnet. Für jede Zeile
$\widetilde a_i^\top$ von $\widetilde A$ wird der Abstand zur zugehörigen
Hyperebene berücksichtigt. Das lineare Programm lautet
\begin{align*}
  \max_{z,r}\quad & r,\\
  \text{unter}\quad
  & \widetilde a_i^\top z
    +\lVert\widetilde a_i\rVert_2r
    \leq\widetilde b_i
    \qquad\text{für alle }i,\\
  & r\geq 0.
\end{align*}
Der optimale Radius $r^*$ beschreibt den Radius einer größtmöglichen Kugel im
relativen Inneren des Polytops. Gilt $r^*>0$, steht ein geeigneter innerer
Punkt für die Eckenberechnung zur Verfügung. Gilt $r^*=0$, reduzieren die
Ungleichungen die Dimension zusätzlich; dieser degenerierte Fall muss separat
behandelt werden.

\subsection{Bestimmung der Ecken}

Das reduzierte Polytop ist durch Halbräume
\[
  \widetilde a_i^\top z\leq\widetilde b_i
\]
gegeben. Ausgehend vom inneren Punkt berechnet die
Halfspace-Intersection-Methode sämtliche Ecken
\[
  z_1,\dots,z_N.
\]
Die Rücktransformation in die ursprünglichen Gewichtskomponenten lautet
\[
  w_i=w_0+Bz_i.
\]

Anschließend werden numerisch identische Ecken zusammengefasst und gegen alle
ursprünglichen Ungleichungen und Gleichungen geprüft. Diese Validierung
verhindert, dass numerisch unzulässige Qhull-Ergebnisse in die
Volumenberechnung gelangen.

\subsection{Volumen aus der konvexen Hülle}

Jedes beschränkte konvexe Polytop ist die konvexe Hülle seiner Ecken:
\[
  P=\conv\{z_1,\dots,z_N\}.
\]
Die konvexe Hülle kann in $d_W$-dimensionale Simplizes zerlegt werden. Ein
Simplex mit den Ecken $v_0,\dots,v_{d_W}$ besitzt das Volumen
\[
  \Vol
  \bigl(
    \conv\{v_0,\dots,v_{d_W}\}
  \bigr)
  =
  \frac{
    \left|
      \det
      \begin{pmatrix}
        v_1-v_0 & \cdots & v_{d_W}-v_0
      \end{pmatrix}
    \right|
  }{d_W!}.
\]
Das Volumen des Polytops ist die Summe der Volumina der Simplizes einer
überlappungsfreien Zerlegung. Die verwendete Convex-Hull-Engine führt diese
geometrische Berechnung intern durch.

Es wird weder ein Raster aufgebaut noch über eine feste Zahl zufälliger Punkte
integriert. Die Berechnung ist daher deterministisch und besitzt keinen
Monte-Carlo-Fehler.

\section{Exakte Wahrscheinlichkeiten für eine Query}

Für eine mögliche Query
\[
  q=(a,b,s)
\]
werden die beiden voll-dimensionalen Antwortregionen definiert als
\begin{align*}
  W_q^{<}(T)
  &=
  W(T)\cap
  \left\{w\mid w_a-s w_b\leq 0\right\},\\
  W_q^{>}(T)
  &=
  W(T)\cap
  \left\{w\mid -w_a+s w_b\leq 0\right\}.
\end{align*}
Die Antwortwahrscheinlichkeiten ergeben sich durch
\begin{align*}
  \Pr(<\mid q,T)
  &=
  \frac{\Vol(W_q^{<}(T))}{\Vol(W(T))},\\
  \Pr(>\mid q,T)
  &=
  \frac{\Vol(W_q^{>}(T))}{\Vol(W(T))}.
\end{align*}

Die Hyperebene
\[
  w_a-s w_b=0
\]
besitzt in einem voll-dimensionalen, kontinuierlich verteilten Gewichtsraum
Volumen null. Deshalb hat eine exakte Gleichheitsantwort unter diesem Modell
Wahrscheinlichkeit null. Wird eine Gleichheit tatsächlich als Nutzereingabe
gegeben, reduziert sie den affinen Raum für alle nachfolgenden Berechnungen;
das Volumen wird danach in diesem neuen intrinsischen Raum bestimmt.

\section{Einbindung in die vorhandene Zielfunktion}

Die exakte Volumenmethode ändert die Zielfunktion des Algorithmus nicht. Sei
$V_h(T)$ der Wert eines Zustands bei noch $h$ verfügbaren Query-Schritten. Im
Terminalfall kann beispielsweise die Anzahl der verbleibenden Kandidaten
verwendet werden:
\[
  V_0(T)=|K(T)|.
\]

Für eine Query $q$ wird der erwartete Folgewert durch die exakten
Antwortwahrscheinlichkeiten gewichtet:
\[
  V_h(T;q)
  =
  \sum_{o\in\{<,>\}}
  \Pr(o\mid q,T)
  V_{h-1}\bigl(T\cup\{(q,o)\}\bigr).
\]
Die gewählte Query ist weiterhin
\[
  q^*
  =
  \arg\min_{q\in Q(T)}V_h(T;q).
\]

Die Volumenberechnung liefert hier ausschließlich
$\Pr(o\mid q,T)$. Die Bedeutung von $V_h$, die Menge $Q(T)$ der zulässigen
Query-Kandidaten und die Minimierung bleiben unverändert.

\section{Algorithmischer Ablauf}

Für die Kandidatenwahrscheinlichkeiten eines Zustands wird folgender Ablauf
verwendet:

\begin{enumerate}
  \item Erzeuge den aktuellen Gewichtsraum $W(T)$ aus Simplexbedingungen und
        beantworteten Queries.
  \item Eliminiere die expliziten Gleichungen und bestimme die Ecken von
        $W(T)$.
  \item Berechne das intrinsische Gesamtvolumen $\Vol(W(T))$.
  \item Für jede Alternative $k$:
    \begin{enumerate}
      \item ergänze alle Optimalitätsbedingungen
            $(x_j-x_k)^\top w\leq 0$,
      \item bestimme die Ecken von $W_k(T)$,
      \item berechne $\Vol(W_k(T))$,
      \item setze
            $p_k(T)=\Vol(W_k(T))/\Vol(W(T))$.
    \end{enumerate}
  \item Prüfe numerisch
        $\sum_k p_k(T)\approx 1$.
\end{enumerate}

Für eine Query-Antwortwahrscheinlichkeit werden statt der
Optimalitätsbedingungen die jeweilige Query-Ungleichung und ihr Gegenstück
eingesetzt. Für die rekursive Query-Bewertung wird dieser Vorgang in den
Kindzuständen wiederholt.

\section{Beispiel mit drei Zielen}

Bei drei Zielen ist der unbeschnittene Gewichtsraum das Dreieck
\[
  W_0
  =
  \conv
  \left\{
    (1,0,0)^\top,
    (0,1,0)^\top,
    (0,0,1)^\top
  \right\}.
\]
Eine Antwort auf die Query
\[
  w_1<2w_2
\]
schneidet das Dreieck mit der Geraden
\[
  w_1=2w_2.
\]
Der verbleibende Gewichtsraum ist wieder ein konvexes Polygon. Die
Optimalitätsbedingungen einer Alternative schneiden dieses Polygon weiter und
erzeugen beispielsweise ein kleineres Dreieck oder Viereck. Die
Kandidatenwahrscheinlichkeit ist dann
\[
  \frac{
    \text{Fläche der Optimalitätsregion}
  }{
    \text{Fläche des aktuellen Gewichtsraums}
  }.
\]
In sieben Zielen ist das Prinzip identisch; an die Stelle von Flächen treten
sechsdimensionale intrinsische Volumina.

\section{Bedeutung von ``exakt''}

Die Methode ist geometrisch und deterministisch exakt:
\begin{itemize}
  \item Es existiert kein Samplingfehler.
  \item Burn-in, Thinning, ESS und R-hat werden nicht benötigt.
  \item Dieselbe Eingabe erzeugt dasselbe Ergebnis.
  \item Die vollständige konvexe Geometrie des Zustands wird ausgewertet.
\end{itemize}

Die Implementierung verwendet jedoch Fließkommaarithmetik und numerische
LP- sowie Convex-Hull-Verfahren. ``Exakt'' bedeutet daher nicht symbolisch
exakt über rationalen Zahlen, sondern numerisch exakt bis auf
Maschinenrundung und Solver-Toleranzen.

Im iterativen Benchmark mit drei, fünf und sieben Zielen, jeweils zehn
Alternativen und bis zu zwanzig beantworteten Queries, wurden 120 Zustände
erfolgreich berechnet. Der größte beobachtete Fehler in der Identität
\[
  \sum_k\Vol(W_k(T))=\Vol(W(T))
\]
betrug relativ
\[
  2{,}85\cdot 10^{-10}.
\]

\section{Laufzeit und Grenzen}

Die Eckenanzahl eines Polytops kann im schlechtesten Fall exponentiell mit der
Dimension und der Anzahl der Halbräume wachsen. Die exakte Methode besitzt
daher keine dimensionsunabhängige Laufzeitgarantie. Für den aktuell
untersuchten Bereich erwies sie sich jedoch als praktikabel.

Im anspruchsvollsten Benchmarkzustand mit sieben Zielen und zwanzig Queries
wurden bis zu 168 Ecken beobachtet. Die Berechnung des Gesamtvolumens und aller
zehn Kandidatenvolumina benötigte im Median etwa $0{,}045$ Sekunden und im
Maximum etwa $0{,}104$ Sekunden auf dem Entwicklungsrechner.

Kritische Fälle sind insbesondere:
\begin{itemize}
  \item sehr hohe Zieldimensionen,
  \item sehr viele Query-Nebenbedingungen,
  \item extrem dünne oder beinahe degenerierte Polytopen,
  \item eine sehr große Zahl separat auszuwertender hypothetischer Queries und
        Kindzustände.
\end{itemize}

Die exakte Methode sollte deshalb durch folgende Prüfungen abgesichert werden:
\begin{itemize}
  \item Validierung aller berechneten Ecken gegen das Nebenbedingungssystem,
  \item Vergleich der Summe der Kandidatenvolumina mit dem Gesamtvolumen,
  \item Erkennung leerer und niedriger-dimensionaler Regionen,
  \item Laufzeit- und Komplexitätsgrenzen für unerwartet große Polytopen.
\end{itemize}

\section{Bezug zur Implementierung}

Die zentralen Bausteine befinden sich in den folgenden Modulen:
\begin{itemize}
  \item \texttt{multistep/src/weight\_space.py}: Aufbau des aktuellen
        Gewichtsraums aus beantworteten Queries,
  \item \texttt{multistep/src/optimality\_region.py}: Ergänzung der
        Optimalitätsbedingungen einer Alternative,
  \item \texttt{multistep/src/polytope\_geometry.py}: Elimination der
        Gleichungen, relativer Chebyshev-Mittelpunkt und Eckenbestimmung,
  \item \texttt{multistep/scripts/benchmark\_iterative\_volume\_methods.py}:
        reproduzierbare Volumenberechnung und iterativer Benchmark.
\end{itemize}

\section{Zusammenfassung}

Die exakte Methode basiert auf fünf Schritten:
\begin{enumerate}
  \item Darstellung des zulässigen Gewichtsraums durch lineare Gleichungen und
        Ungleichungen,
  \item Ergänzung von Query- oder Optimalitätsbedingungen,
  \item Projektion in den intrinsischen affinen Raum,
  \item Bestimmung der Ecken und des Volumens der konvexen Hülle,
  \item Berechnung der Wahrscheinlichkeiten als Volumenquotienten.
\end{enumerate}

Damit können die für die vorhandene Zielfunktion benötigten
Kandidaten- und Query-Antwortwahrscheinlichkeiten ohne Monte-Carlo-Sampling
bestimmt werden.

\end{document}
